% Figure source: positional rotations in RoPE.
% Uses the standard LaTeX picture environment; no additional package is needed.

\begin{figure}[htbp]
    \centering
    \setlength{\unitlength}{1mm}
    \resizebox{\textwidth}{!}{%
    \begin{picture}(158,66)
        \put(0,53){\framebox(28,10){Query $\mathbf{q}_m$}}
        \put(33,58){\vector(1,0){9}}
        \put(42,53){\framebox(36,10){$R_m=R(m\boldsymbol{\theta})$}}
        \put(80,58){\vector(1,0){9}}
        \put(89,53){\framebox(34,10){$\widetilde{\mathbf{q}}_m$}}

        \put(0,31){\framebox(28,10){Key $\mathbf{k}_n$}}
        \put(33,36){\vector(1,0){9}}
        \put(42,31){\framebox(36,10){$R_n=R(n\boldsymbol{\theta})$}}
        \put(80,36){\vector(1,0){9}}
        \put(89,31){\framebox(34,10){$\widetilde{\mathbf{k}}_n$}}

        \put(0,10){\makebox(0,0)[l]{$
        \widetilde{\mathbf{q}}_m^{\mathsf T}
        \widetilde{\mathbf{k}}_n
        =
        \mathbf{q}_m^{\mathsf T}R_{n-m}\mathbf{k}_n
        $}}
    \end{picture}%
    }
    \caption{Rotary Position Embedding applies position-dependent rotations
    to queries and keys. Their inner product contains a rotation determined
    by the relative displacement between positions. Schematic based on
    \textcite{zhang2024foundmiddle}.}
    \label{fig:pe-rope-relative-position}
\end{figure}
